\chapter{选择计算参数}\label{ch:theory}

完成一次计算之后，接下来的问题通常是：这些参数为什么这样选，换一个体系时应该改哪些，以及怎样判断精度已经足够。本章围绕界面中的实际设置介绍电子结构计算的基本概念。阅读时可以保留第\ref{ch:first}章的氢分子项目，一边查看参数，一边理解它们所描述的物理问题。

一套完整的计算条件包括几何结构、电子模型和数值精度。几何结构给出原子种类、位置与晶格；电子模型给出泛函、赝势和自旋等假设；数值精度则涉及截断能、k 点采样和迭代收敛条件。三者各有作用。数值上充分收敛的计算，仍然是在所选模型下求得的答案；模型选得合适，也需要足够的数值精度才能得到稳定结果。

\section{从原子位置到电子密度}

\subsection{密度泛函理论解决什么问题}
\index{Midu fanhan lilun@密度泛函理论}

在给定原子位置时，电子并不是静止在原子周围的一组小球。它们的运动和相互作用共同决定体系的能量、化学键和许多材料性质。直接追踪多电子波函数十分困难，密度泛函理论改用空间中的电子密度来描述基态问题。

电子密度可以理解为电子在空间中分布的程度。密度较大的区域通常有较多电子，但它并不等于查看器中原子球的大小。球棍模型帮助我们辨认几何结构，电子密度则来自计算；两者的图像形态和用途不同。

KSSOLV 通过 Kohn--Sham 方程求解电子态。用于求解的有效势取决于电子密度，而电子密度又由电子态产生，因此不能只计算一次便结束。求解器需要不断更新密度，直到方程所用的密度与求得的密度互相一致，这就是自洽。


在本书非磁性、PBE赝势教程中，可以把所求的方程写为
\begin{equation}
  \begin{aligned}
  \widehat H[\rho]\psi_{n\mathbf k}
    &=\varepsilon_{n\mathbf k}\psi_{n\mathbf k},\\
  \widehat H[\rho]&= -\frac12\nabla^2
  +V_{\mathrm{loc}}^{\mathrm{ps}}+\widehat V_{\mathrm{NL}}^{\mathrm{ps}}
  +V_{\mathrm H}[\rho]+V_{\mathrm{xc}}[\rho].
  \end{aligned}
  \label{eq:ks-operator}
\end{equation}
这里采用原子单位：长度为Bohr，能量为Ha。动能项描述电子的空间变化；局域和非局域赝势描述原子实；Hartree势与交换关联势随密度更新。界面中的\ui{Pseudopotential}和\ui{Functional}分别决定相应部分。HSE06还含非局域交换处理，不能仅用上式中的局域交换关联势概括。

非磁性设置\ui{Nspin}=1时，电子密度由这些态加权得到：
\begin{equation}
  \begin{aligned}
  \rho(\mathbf r)&=2\sum_{\mathbf k}w_{\mathbf k}
       \sum_n f_{n\mathbf k}|\psi_{n\mathbf k}(\mathbf r)|^2,\\
  N_{\mathrm e}&=\int_{\Omega}\rho(\mathbf r)\,\mathrm d^3r.
  \end{aligned}
  \label{eq:ks-density}
\end{equation}
波函数在一个计算晶胞内归一化。前面的2计入非磁性自旋简并；满占据的一个空间态取$f=1$，对应两个电子，不应再把$f$写成2。分子教程只有一个$\Gamma$点，其权重为1。赝势计算中的$N_{\mathrm e}$是实际处理的价电子数，密度使用电子数每单位体积表示；它与含负电子电荷的电荷密度有符号上的区别。

\begin{table}[htbp]
\centering
\caption{阅读电子结构公式时的常用符号。}
\begin{tabularx}{\linewidth}{lX}
\toprule
符号 & 含义及与设置的联系\\
\midrule
$\psi_{n\mathbf k}$、$\varepsilon_{n\mathbf k}$ & 第$n$个电子态及其能量；$\mathbf k$标记晶体中的波矢\\
$\rho$、$\Omega$ & 价电子密度及计算晶胞的体积；原子单位下分别用$\mathrm{Bohr}^{-3}$和$\mathrm{Bohr}^{3}$\\
$f_{n\mathbf k}$ & 占据权重，由电子数、\ui{Smearing}和\ui{Temperature}共同确定\\
$w_{\mathbf k}$ & $k$点积分权重；一个自旋通道内的权重和为1\\
\bottomrule
\end{tabularx}
\end{table}

改变\ui{Extra bands}可以增加参与求解的空态，但不会改变式\eqref{eq:ks-density}所要求的电子数。选择自旋分辨或非共线模式时，需要分别处理相应通道，不能直接沿用这里的简并因子2。

\subsection{为什么固定原子也要迭代}

可以把一次 SCF 看成反复改进一份电子分布的过程。首先建立一个初始密度，用它构造有效势；随后求解电子态，根据这些态和电子占据数得到新密度；再把新旧信息适当混合，作为下一步的起点。

\begin{figure}[htbp]
\centering
\begin{tikzpicture}[x=1mm,y=1mm,font=\sffamily\small,
  stage/.style={draw=teal,fill=paleblue,rounded corners=1.5pt,
    minimum width=31mm,minimum height=11mm,align=center,text=brand},
  flow/.style={->,>=stealth,line width=.7pt,draw=teal}]
  \node[stage] (input) at (15,32) {输入密度\\$\rho_{\mathrm{in}}$};
  \node[stage] (potential) at (55,32) {构造有效势};
  \node[stage] (states) at (95,32) {求电子态与占据};
  \node[stage] (output) at (95,11) {得到新密度\\$\rho_{\mathrm{out}}$};
  \node[stage] (check) at (55,11) {达到收敛条件？};
  \node[stage] (mix) at (15,11) {混合新旧密度};
  \node[text=brand,align=center] (done) at (55,-6) {保留收敛结果};
  \draw[flow] (input) -- (potential);
  \draw[flow] (potential) -- (states);
  \draw[flow] (states) -- (output);
  \draw[flow] (output) -- (check);
  \draw[flow] (check) -- node[above,text=brand] {否} (mix);
  \draw[flow] (mix) -- (input);
  \draw[flow] (check) -- node[right,text=brand] {是} (done);
\end{tikzpicture}
\caption{电子自洽循环（示意）。在原子位置不变的条件下更新电子密度，直到输入与输出满足收敛要求。}
\label{fig:scf-cycle}
\end{figure}


这个过程中，原子位置可以完全不变。因此，SCF 的“第十步”不表示原子移动了十次。结构优化还包含另一层循环：先在一个几何构型上完成电子自洽，再根据原子力移动原子，然后对新构型重新计算电子结构。区分电子迭代和原子位置迭代，是阅读日志的重要基础。

\subsection{自洽与非自洽计算的分工}
\index{Ziqiachang@自洽场（SCF）}

SCF 的目标是建立与结构和模型一致的电子密度。能带和态密度还需要在更多或不同的 k 点上求电子能量，这时可以使用已经收敛的密度，保持有效势不再随这些附加采样反复自洽。这类计算通常称为非自洽计算，简称 NSCF。

在工作流中，\ui{SCF} 先提供收敛密度；\ui{NSCF} 指定后续求解所需的算法选项；\ui{Band Calculation} 或 \ui{DOS} 再根据相应的路径或网格执行计算。能带路径和 DOS 网格描述的是不同的采样任务，不应因为它们都含有 k 点，就把两者设成同一个列表。

\section{选择分子还是晶体}

\subsection{计算对象的含义}

\ui{Build Molecule} 中的 \ui{Molecule} 与 \ui{Crystal} 决定如何组织计算对象。对于孤立分子教程，选择 Molecule，并为分子保留合适的模拟盒；对于周期性材料，选择 Crystal，并给出晶格和 k 点采样。

“分子”并不意味着计算完全没有盒子。平面波方法仍然需要一个空间区域来展开波函数。盒内分子与其周期重复像之间的相互作用，应通过足够的真空距离减小。对于带电、强极性或空间分布很宽的体系，所需盒尺寸可能与中性紧凑分子不同。

分子文件、晶体文件与计算类型也不能完全等同。CIF 通常保存晶体，但也可以保存分子示例的模拟盒和其中的分子坐标。XYZ 通常保存分子坐标，导入后软件可以为它构造一个盒子。最终仍应检查实际几何信息，并在节点中明确选择计算类型。

\subsection{晶格向量与原子坐标}

周期结构由三个晶格向量和晶胞内的原子位置组成。晶格的长度、夹角决定晶胞形状；原子位置可以用直角坐标表示，也可以用沿三个晶格向量的分数坐标表示。

分数坐标 $(0.25,0.25,0.25)$ 表示沿三个晶格向量各取四分之一后相加，并不总等于三个直角坐标都为晶格常数的四分之一。只有在相应的正交晶胞中，二者才有这种简单关系。对于六方、单斜等晶胞，输入坐标时尤其应确认当前坐标类型。


用$\mathbf a_i$表示三个实空间晶格向量，$f_i$表示分数坐标，则位置与倒格向量满足
\begin{equation}
  \begin{gathered}
  \mathbf r=\sum_{i=1}^{3}f_i\mathbf a_i,
  \qquad \mathbf k=\sum_{j=1}^{3}\kappa_j\mathbf b_j,\\
  \mathbf a_i\!\cdot\!\mathbf b_j=2\pi\delta_{ij}.
  \end{gathered}
  \label{eq:reciprocal-convention}
\end{equation}
$\delta_{ij}$在$i=j$时为1，否则为0。这里的倒格约定包含$2\pi$；$\kappa_j$是倒格基底下的无量纲坐标，而$\mathbf k$本身具有长度倒数的单位。因而路径中的$(0.5,0,0)$表示沿第一倒格向量取一半，并非实空间中的半个晶格长度。

一个晶体可以使用原胞、常规晶胞或超胞表示。它们可能描述同一种无限周期结构，却包含不同数量的原子，也有不同大小的布里渊区。超胞越大，单位晶胞内的电子态越多；同样的整数 k 点网格不再代表相同的采样密度。

\subsection{单位的三个层次}

结构编辑和常用文件格式通常使用 Å，数值计算中的几何量使用 Bohr，电子能量则可能在不同位置以 Ha、Ry 或 eV 表示。软件会在导入和求解之间完成必要的转换，用户仍需在读取外部数据或整理结果时关注单位。

\begin{table}[htbp]
\centering
\caption{本书常见量与单位}
\begin{tabularx}{\linewidth}{lX}
\toprule
量 & 常见表示与阅读位置\\
\midrule
几何长度 & 查看器、CIF 和通常的 XYZ 用 Å；数值接口可使用 Bohr\\
截断能 & Build Molecule 的 Ecut 用 Hartree，即 Ha\\
总能量 & 收敛曲线左轴用 Ha；日志也可能同时列出 Ry\\
能带与 DOS 能量 & eV，本书教程以费米能为参考零点\\
电子温度 & Temperature 用 K\\
原子力 & 数值计算记录采用原子单位，比较外部资料时应换算\\
\bottomrule
\end{tabularx}
\end{table}

常用换算为 $1\,\mathrm{Ha}=2\,\mathrm{Ry}\approx27.2114\,\mathrm{eV}$，$1\,\mathrm{Bohr}\approx0.529177\,\text{Å}$。如果只是检查输入，使用这些近似值即可；批量整理数值时，应使用程序中的转换常数，避免在不同表格中混用过多种舍入值。

\section{泛函：选择电子相互作用的近似}

\subsection{常用泛函}

\ui{Functional} 提供 PBE、PZ 和 HSE06。它们采用不同方式近似电子的交换关联能。PZ 属于局域密度近似一类，PBE 属于广义梯度近似，HSE06 属于屏蔽杂化泛函。它们的区别影响计算得到的平衡结构、相对能量和电子能级。

对于学习工作流和一般的初步材料计算，PBE 是便于开始的选择。本书的基础教程统一使用 PBE，使读者能够把注意力放在结构和计算过程上。切换到其他泛函时，不能只比较运行时间，还应考虑赝势是否合适、研究目标是什么，以及需要比较的文献采用了哪些条件。

HSE06 引入了一部分精确交换，计算量通常高于半局域泛函。它还需要相应的交换算法和后续计算条件。初次学习时，先完成 PBE 的 SCF、能带和 DOS，再用相同的小体系学习杂化泛函的算法设置。

\subsection{怎样比较不同泛函}

比较泛函时，应先决定比较的是“同一结构上的电子模型差异”，还是“各自平衡结构上的结果差异”。前者保持几何不变，只改变电子模型；后者允许不同模型得到不同的最稳定结构。这两种比较回答的问题不同。

例如，在固定键长上比较两个氢分子总能量，并不能直接判断哪种泛函给出的平衡键长更接近实验。若研究目标是键长，应该分别优化结构；若目标是固定结构下的电子能级，则应使用相同几何。写计算说明时把这个选择说明白，结果会更容易理解。

\section{赝势：怎样描述原子实}
\index{Yanshi@赝势}

赝势把紧密束缚的内层电子和原子核的影响合并为有效描述，使计算主要处理价电子。这是平面波计算能够高效处理多种元素的重要原因。赝势本身包含建模选择，因此它也是计算条件的一部分。

\ui{Pseudopotential} 的 \ui{default} 使用随软件提供的默认赝势集合。本书的氢和硅教程使用内置 PBE-1.3 ONCV 文件。在运行消息中可以查看实际加载的文件名称，确认每一种元素都找到了对应的赝势。

泛函和赝势应配合使用。更改 Functional 不表示软件会自动选择一套新的赝势，也不表示所有组合都同样合适。开始新的材料研究时，应记录赝势来源、版本和价电子选择，并在同一组可比较的计算中保持一致。

\subsection{为什么不同赝势需要不同截断能}
\index{Jieduanneng@截断能}

有的赝势产生较平滑的价电子波函数，较少的平面波便能描述；有的需要更高的截断能。截断能是否足够，取决于体系和所用赝势，而不是只取决于元素的原子序数或原子数量。

当一个结构含有多种元素时，应针对完整体系进行检查。某个元素对平面波精度要求较高，往往会决定整个计算需要的截断能。氧化物、过渡金属或含氢体系应分别进行截断能收敛检查。

\section{平面波截断能}
\index{Pingmianbo@平面波}

\subsection{\headinglatin{Ecut} 控制什么}

平面波把电子波函数表示为空间波的叠加。波长越短，能够描述越细的空间变化，相应的动能越高。截断能规定保留哪些平面波，从而控制表示电子态的精细程度。


晶体中的一个电子态可以展开为
\begin{equation}
  \begin{aligned}
  \psi_{n\mathbf k}(\mathbf r)
   &=\frac{1}{\sqrt{\Omega}}\sum_{\mathbf G}
      c_{n\mathbf k}(\mathbf G)
      \exp\!\left[\mathrm i(\mathbf k+\mathbf G)\cdot\mathbf r\right],\\
  T_{\mathbf k+\mathbf G}&=\frac{|\mathbf k+\mathbf G|^2}{2}.
  \end{aligned}
  \label{eq:plane-wave-expansion}
\end{equation}
$\mathbf G$是倒格向量的整数线性组合，$c$是待求系数；采用Bohr及其倒数时，右侧动能的单位为Ha。对$\Gamma$点，$\mathbf k=0$，基础平面波集合满足$|\mathbf G|^2/2\leq E_{\mathrm{cut}}$。提高截断能，就是把更远处、波长更短的分量纳入表示。

\begin{figure}[htbp]
\centering
\begin{tikzpicture}[x=1mm,y=1mm,font=\sffamily\small]
  \node[text=brand] at (26,47) {$\Gamma$点的基础平面波集合};
  \node[text=brand] at (88,47) {波长与空间细节};
  \fill[paleblue] (26,22) circle (18);
  \foreach \x in {-24,-18,...,24}
    \foreach \y in {-18,-12,...,18} {
      \pgfmathtruncatemacro{\inside}{\x*\x+\y*\y<=324}
      \ifnum\inside=1
        \fill[brand] ({26+\x},{22+\y}) circle (.75);
      \else
        \fill[muted!35] ({26+\x},{22+\y}) circle (.55);
      \fi
    }
  \draw[teal,line width=.8pt] (26,22) circle (18);
  \draw[->,draw=muted,line width=.4pt] (1,22) -- (53,22);
  \draw[->,draw=muted,line width=.4pt] (26,0) -- (26,43);
  \node[anchor=west,text=muted] at (51,19) {$G_x$};
  \node[anchor=west,text=muted] at (27,42) {$G_y$};
  \node[fill=paleblue,inner sep=1pt,text=brand] at (26,22) {$\Gamma$};
  \draw[draw=muted!40,line width=.4pt] (66,31) -- (112,31);
  \draw[draw=muted!40,line width=.4pt] (66,12) -- (112,12);
  \draw[teal,line width=.8pt,domain=0:44,samples=100]
    plot ({67+\x},{31+3*sin(\x*16.3636)});
  \draw[brand,line width=.8pt,domain=0:44,samples=140]
    plot ({67+\x},{12+3*sin(\x*49.0909)});
  \node[text=ink] at (89,39) {较长波长};
  \node[text=ink] at (89,20) {较短波长};
  \node[text=ink] at (26,-7) {提高Ecut，保留更多分量};
  \node[text=ink] at (89,-7) {短波分量需要更高动能};
\end{tikzpicture}
\caption{平面波截断的二维示意。左图蓝点为截断圆内保留的倒格分量；实际计算使用三维集合。右图说明波长与空间变化的关系。}
\label{fig:plane-wave-cutoff-concept}
\end{figure}


非零$k$点的动能取决于$\mathbf k+\mathbf G$。晶体计算可以共用按基础$\mathbf G$选择的集合，因此图中的圆应理解为$\Gamma$点示意，而不是每个$k$点都重新裁切出完全相同的动能球。实际使用时仍通过收敛检查判断\ui{Ecut}是否足够。

提高 \ui{Ecut (Hartree)} 通常会增加平面波数、内存需求和每一步计算时间。它改善的是基组描述，并不会自动增加 k 点采样，也不会自动放宽或收紧电子自洽条件。几个精度参数之间相互影响，但不能互相替代。

\subsection{怎样做截断能收敛检查}

先固定结构、盒子、泛函、赝势和 k 点，再依次提高 Ecut。例如，为一个小型示例体系计算 $15$、$20$、$25$、$30$ Ha 四个值，记录每次收敛后的能量和所关心的性质。这里的序列只是检查方法的示例，最终范围应根据实际体系调整。

判断时关注相邻设置之间的变化，而不是某个总能量看起来是否“足够负”。如果研究的是相对能量，可以记录每原子能量差或两个结构之间的能量差；如果研究的是几何优化，应同时检查力和优化后的尺寸。

\begin{steps}
\item 从同一模板另建一条工作流，或先保存已有工作流为 \file{.wf}，再通过从文件导入建立副本，保留相同结构。
\item 每次只修改 Ecut，并让 SCF 达到同样的收敛条件。
\item 把结果整理为截断能与目标性质的表格。
\item 选择在所需精度范围内已经稳定的设置，再进行正式计算。
\end{steps}

SCF 未收敛的结果不宜混入这样的比较。否则观察到的差异可能来自迭代尚未结束，而不是真正的基组精度变化。收敛检查通常比盲目选择很大的参数更节省整体工作量。

\section{k 点与布里渊区采样}
\index{K dian caiyang@k 点采样}

\subsection{为什么晶体需要 k 点}

周期结构中的电子态可以用波矢 k 来标记。许多总能量和密度有关的量需要对整个布里渊区求和或积分，实际计算则用有限数量的 k 点近似。\mbox{\ui{Autokpts}} 指定规则网格，前三个整数给出三个方向上的网格数，后三项指定相应的网格位移。


对布里渊区内某个量$Q(\mathbf k)$，区域平均可用加权和近似：
\begin{equation}
  \langle Q\rangle_{\mathrm{BZ}}
      \approx\sum_{\mathbf k}w_{\mathbf k}Q(\mathbf k),
  \qquad \sum_{\mathbf k}w_{\mathbf k}=1.
  \label{eq:k-point-average}
\end{equation}
规则完整网格可使用等权重。增加网格密度，改善的是这类区域采样；沿一条路径加入更多点，只能使路径上的曲线更细，不能代替整个布里渊区的积分网格。

例如 \mbox{\texttt{[2 2 2 0 0 0]}} 描述本书硅教程中的起始网格。它是方便学习的小网格，不是硅的所有性质都已充分收敛的结论。增加到更密网格后，可以检查总能量、力或能带边缘是否已经稳定。

对于细长或扁平的晶胞，不同方向不一定需要相同的整数。实空间周期越长，相应倒空间方向的区间越短，所需采样数通常也会改变。应结合晶胞形状和目标精度选择网格，而不是只追求三个数字相同。

\subsection{为任务选择采样方式}

SCF 网格用于建立电子密度和有效势。能带路径选择若干具有对称意义的方向，展示能量如何随 k 变化。DOS 网格则覆盖整个布里渊区，用于统计一定能量附近有多少电子态。

\begin{figure}[htbp]
\centering
\begin{tikzpicture}[x=1mm,y=1mm,font=\sffamily\small]
  \node[text=brand] at (21,43) {SCF 与 DOS 网格};
  \node[text=brand] at (85,43) {能带路径};
  \fill[paleblue] (3,0) rectangle (39,36);
  \fill[paleblue] (67,0) rectangle (103,36);
  \foreach \x in {3,12,21,30,39}
    \draw[teal!30,line width=.3pt] (\x,0) -- (\x,36);
  \foreach \y in {0,9,18,27,36}
    \draw[teal!30,line width=.3pt] (3,\y) -- (39,\y);
  \draw[teal,line width=.65pt] (3,0) rectangle (39,36);
  \draw[teal,line width=.65pt] (67,0) rectangle (103,36);
  \foreach \x in {7.5,16.5,25.5,34.5}
    \foreach \y in {4.5,13.5,22.5,31.5}
      \fill[brand] (\x,\y) circle (1.1);
  \draw[brand,line width=1pt] (85,18) -- (103,18) -- (103,36) -- cycle;
  \foreach \x/\y in {85/18,88/18,91/18,94/18,97/18,100/18,103/18,
    103/21,103/24,103/27,103/30,103/33,103/36,
    100/33,97/30,94/27,91/24,88/21}
      \fill[brand] (\x,\y) circle (.85);
  \node[fill=paleblue,inner sep=1pt,text=brand,anchor=north east] at (84,17) {$\Gamma$};
  \node[text=brand,anchor=west] at (105,18) {X};
  \node[text=brand,anchor=south west] at (104,35) {M};
  \node[text=ink] at (21,-7) {覆盖整个区域};
  \node[text=ink] at (85,-7) {沿选定方向采样};
\end{tikzpicture}
\caption{$k$点网格与能带路径的区别（二维示意）。前者用于区域积分或统计取样，后者展示选定路径上的变化。图中以正方形区域和三角形路径说明两种取样方式。}
\label{fig:k-sampling}
\end{figure}


这三类采样可以在同一项目中并存。例如，一条周期计算流程可以先以 $2\times2\times2$ 网格演示 SCF，再沿推荐高对称路径计算能带，或以 $4\times4\times4$ 网格计算总态密度。后续任务继承所属流程的自洽密度，使用各自的采样方式。本书的硅 DOS 入门流程使用八原子晶胞，能带流程使用二原子标准原胞，两者各自完成 SCF。第\ref{ch:results}章还介绍二原子原胞项目中保存的 DOS 数值与重绘方法；阅读时应核对数据所属的晶胞。

能带图上的一条平滑曲线并不能说明整个布里渊区积分已经准确。相反，一张很密的 DOS 网格也不会让一条仅含两个端点的能带路径自动变得平滑。选择参数时应先问清楚当前节点要完成哪一类任务。

\subsection{高对称路径的含义}

对称性分析根据晶格和原子排列推荐高对称点与路径。路径上的标签，如 $\Gamma$、X、L 等，指向倒空间中特定的位置；相同标签在不同晶系中并不总有相同坐标。

改变晶胞表示、原子排列或对称性容差，可能影响推荐路径。比较两张能带图时，除了泛函和能量零点，还应检查它们是否采用相同晶胞和同一路径。不能只按横轴标签相似就把曲线逐点相减。

本书使用 \ui{Symmetry \& K-Path} 生成路径，在 \ui{Band Calculation} 中用 \mbox{\ui{Segment Points}} 设置每段采样点数。这一数值包含端点。初次运行可以用较少点数观察整体结构；需要平滑图形或精细比较曲率时，再增加采样。

\section{设置电子占据}

\subsection{电子占据与费米能}
\index{Feimineng@费米能}

计算得到许多电子态后，还需要决定有多少电子占据各个态。费米能与电子总数和占据方式有关，是阅读能带和 DOS 时常用的参考。金属中费米能附近有部分占据的态，绝缘体和半导体则在占据态与未占据态之间可能存在能隙。

数值计算中的费米能不是各个体系都可以直接比较的绝对真空能级。本书能带图把费米能移到零点，便于辨认其附近的电子态。如果要比较功函数、能带对齐或不同体系之间的绝对能级，需要进一步建立一致的能量参照。

\subsection{展宽与电子温度}
\index{Zhankuan@展宽与电子占据}

\ui{Smearing} 提供 fermi-dirac、cold、gaussian 和 mp。展宽方法使占据在能量上较平滑，常用于改善金属或接近简并态体系的布里渊区积分和收敛。\ui{Temperature (K)} 给出相应的电子温度参数。

电子温度参数不等于给原子加热，也不意味着计算已经包含某一温度下的原子振动。一个在 $300$ K 电子占据下运行的静态 SCF，原子仍然可以停在固定位置。需要有限温度结构性质时，应选择与研究问题相符的方法。

改变展宽和温度会影响占据及相应的能量量。比较不同结果时，应说明使用的是总能量、自由能还是其他经过处理的量，并保持定义一致。对于初学教程，先按给定条件运行，再单独研究展宽变化，阅读结果会清楚得多。

\subsection{额外能带}
\index{Ewai nengdai@额外能带}

电子数决定至少需要多少占据态，但能带和 DOS 往往还要查看费米能以上的空态。\ui{Extra bands} 增加求解的未占据能带数量。它不是给体系增加电子，也不会直接改变原子的价态。

若图中只有占据态而缺少所关心的导带区域，应检查额外能带数，而不是只把纵轴范围拉得更大。坐标轴只能显示已有数据，不能生成没有计算的电子态。本书硅教程设置额外 $4$ 条能带，作为观察低能空态的起点；更高能区需要相应增加。

\section{自旋与磁性}
\index{Zixuan@自旋}

\ui{Nspin} 提供 $1$、$2$、$4$ 等设置，分别对应不同的自旋处理。本书的非磁性示例使用 $1$。存在未配对电子、磁有序或自旋轨道耦合时，应根据实际物理问题选择适当模式，并检查相关高级设置。

自旋设置不是一个仅影响速度的开关。氢原子与氢分子具有不同的电子占据情况；若计算分子解离能，原子参考态的自旋处理会直接影响结果。过渡金属体系还可能存在多个自洽解，初始磁性和几何结构都会影响最后得到的状态。

本书的基础案例不使用磁性、非共线或自旋轨道耦合。学习时可以先把非磁性工作流掌握，再为一个明确的磁性问题建立新项目。这样更容易区分结构、采样和自旋各自造成的变化。

\section{让自洽计算收敛}

\subsection{对角化方法}

\ui{DiagMethod} 控制电子态的数值求解方法。\ui{davidson\_qe}、\ui{lobpcg} 等方法针对同一电子结构问题，采用不同的迭代策略。充分收敛时，它们应给出一致的物理结果，主要区别在计算效率和求解过程。

比较算法时，应使用相同输入和目标精度，以各自收敛后的结果判断一致性。本书的 SCF 例使用默认方法，后续 NSCF 例保留其默认求解设置。

\subsection{最大步数与收敛精度}

\ui{Max SCF steps} 是允许的最大迭代次数。把它从 $100$ 改为 $200$，只给求解器更多机会，不会自动把精度提高一倍。真正的收敛判断由误差和相应阈值决定。到达最大步数但未满足收敛条件，应视为需要继续处理的计算。

基础 SCF 面板提供方法、最大步数和力等选项。密度阈值与混合参数属于进一步的数值设置，脚本接口可以使用 \texttt{options.SCF} 配置。本书的基础教程沿用默认密度阈值 $10^{-8}$。需要调整密度阈值与混合参数时，可在完整MATLAB\regmark{}中使用这一脚本接口。


设一次电子迭代的输入与输出密度之差为
\begin{equation}
  \Delta\rho=\rho_{\mathrm{out}}-\rho_{\mathrm{in}},
  \qquad e_{\mathrm{SCF}}<\varepsilon_{\mathrm{SCF}}.
  \label{eq:scf-stopping}
\end{equation}
左式说明比较的是哪两份密度；右式说明所用误差指标必须低于阈值。默认密度混合计算把$\Delta\rho$在倒空间中作加权比较，得到$e_{\mathrm{SCF}}$，并与默认的$10^{-8}$比较。日志虽使用\file{Rel Rho Err}这一名称，其值不能直接读成百分比，也不是两份密度之差的范数除以输入密度的范数。判断本书实例是否收敛，应将这一指标与对应阈值比较，并查看最终状态。

\subsection{密度混合的作用}
\index{Midu hunhe@密度混合}

一次求解得到的新密度未必是下一次最好的起点。密度混合把当前结果和历史信息组合起来，减小来回振荡并加快收敛。不同体系对混合的响应不同，金属、大晶胞或磁性体系可能需要更仔细的设置。

若误差持续在两个水平之间摆动，可以先检查结构和占据，再考虑混合和采样。增加最大步数可以让一个已经稳定下降的序列继续走完，却未必能解决长期振荡。处理这类问题时，保留完整日志比只保存最后一行报错更有帮助。

\section{几何优化中的两个精度}
\index{Jihe youhua@几何优化}

几何优化要求每一个候选结构都先完成电子计算，然后根据力更新位置。电子密度收敛精度决定每一步能量和力的可靠程度，几何优化阈值则决定何时认为原子位置已经足够稳定。


原子$I$的力是总能量对该原子位置的负梯度：
\begin{equation}
  \mathbf F_I=-\nabla_{\mathbf R_I}E,
  \qquad F_{I\alpha}=-\frac{\partial E}{\partial R_{I\alpha}}
  \quad(\alpha=x,y,z).
  \label{eq:atomic-force}
\end{equation}
在计算采用的原子单位中，力以Ha/Bohr表示。力指向能量下降的局部方向；总能量已经变化很小，并不自动意味着每个原子的力都很小。优化记录还可能列出最大原子力、均方根力或整体梯度范数，比较时应先确认列名与定义。

\ui{RelaxTolerance} 与 \ui{MaxRelaxSteps} 分别控制几何优化的停止尺度和最大工作量。本书氢分子优化例明确输入 $10^{-4}$ 和 $100$。图形面板和脚本接口的默认值可以不同，记录时应写出实际采用的条件。

最终结果应同时检查几何优化是否收敛、最后一次电子计算是否收敛，以及最后的原子力是否符合研究所需。只看到键长不再明显变化，仍不足以判断结构已经达到稳定位置。

\section{选择杂化交换算法}

在Build Molecule的\ui{Functional}中选择HSE06后，计算会包含杂化泛函的交换处理。SCF节点高级选项中的\ui{ExxMethod}用于选择这一处理的算法，可见选项有default、normal、ace、qrcp和kmeans。泛函决定要使用的近似，交换算法决定怎样组织其中的计算；它们不是可以互相替代的两个下拉框。

初次使用时先保留default。需要比较算法时，在一个已经有基态计算记录的小体系上另存项目，只改变ExxMethod，核对总能量、收敛状态和用时。比较期间保持泛函、模型尺寸和采样一致。使用PBE等半局域泛函时，保持default即可。

\section{区分基态与激发态}

SCF的轨道本征值、两个轨道的能量差和多电子体系的激发能是不同的量。基态计算已经收敛，只说明该基态问题得到了解答；希望研究激发时，还要选择响应方法、占据轨道和空轨道组成的计算空间。

TDDFT任务中的\ui{nvbands}和\ui{ncbands}分别是参与响应计算的占据与空轨道数，\ui{nroots}是要求的激发态数。它们都不是原子数，也不是k点数。首先在Build Molecule中求出足够的额外能带，再在TDDFT里选择其中的空轨道。以一个占据轨道和四个空轨道为例，最多只有四个这样的轨道对，不能据此要求十个激发根。

\ui{Use TDA}使用Tamm--Dancoff近似；\ui{Use IPA}和\ui{Use RPA}用于不同响应近似的比较。学习时先使用一套明确的设置，再单独改变近似观察差异。TDDFT的输入准备、适用范围和参数示例见第\ref{ch:workflows}章\ref{sec:tddft-workflow}节。吸收峰需由相应的响应计算得到。


\section{为一个新体系选择参数}

开始新项目时，先明确研究目标，例如比较吸附位置的相对能量、观察应变对能带的影响，或寻找分子的平衡结构。建立能够运行的小规模设置，核对结构、元素和电子数，再围绕目标物理量逐项提高精度。保留一份基准结果，每次改变一个条件，后续比较才有清楚的依据。

随项目保存简短的条件表，记录结构来源与处理、晶胞、泛函、赝势、截断能、k点、占据、自旋及两类收敛条件。报告列出影响结论的主要信息，完整工作流和输入文件留在项目资料中，便于以后理解和重复计算。
